
% ============================================================================
\clearpage
\section*{v4.8.0-r04 Source-Only: Decoupling from the Riemann Source and Standalone Sphaera Transformations}
\addcontentsline{toc}{section}{v4.8.0-r04 Source-Only: Decoupling from the Riemann Source and Standalone Sphaera Transformations}
\label{sec:v480-r04-riemann-decoupling}
% ============================================================================

\paragraph{Purpose.}
The Companion cites the Riemann source \cite{brendecke2026rh} at 23 sites for seven distinct
results.  A full audit of that source (v4.15) and a recomputation of its central formulas
show that its proof chain for RH does not hold.  The author used the source as a starting
point; what the Companion needs from it are the Heisenberg compensator and the Sphaera model
of the basic transformations (reflection, cascade, triole hierarchy).  This layer records the
audit, gives standalone definitions of those three objects that do not depend on any RH
statement, repairs the triole hierarchy where it fails, and reclassifies every other
dependency.  Earlier layers are not rewritten.

\noindent\textbf{Source-only status.}
This v4.8.0-r04 layer is source-only.  It does not build a PDF, does not assign a new DOI,
and does not authorize a public release.  The v4.8.0 Version DOI
\href{https://doi.org/10.5281/zenodo.23365244}{10.5281/zenodo.23365244} and the Concept DOI
\href{https://doi.org/10.5281/zenodo.19210728}{10.5281/zenodo.19210728} are unchanged by it.

\subsection*{v4.8.0-r04 audit of the Riemann source}
\label{subsec:v480-r04-source-audit}

The combinatorial orbit theory (L1--L11, Axioms W and G, K1), the arithmetic identification
\(\mathcal Z=\zeta\) for \(\operatorname{Re}s>1\) and the functional-equation sketch are correct.
The following findings concern the steps that carry the claim.  Numerical checks are recorded
in the audit scripts of the working folder \texttt{llm/}.

\paragraph{A1 --- explicit formula.}
The source states \(\Theta(\tau)=1-\Sigma(\tau)+R(\tau)\) with
\(\Theta=\sum_{p,m}m^{-1}\delta(\tau-m\log p)\).  Against a smooth test function supported in
\((0.6,3.2)\) with 300 zeros the two sides give \(2.1587\) and \(0.5198\).  The correct form,
which matches to six digits (\(4.393268\) on both sides), is
\[
\tau\,\Theta(\tau)=\sum_{n}\Lambda(n)\,\delta(\tau-\log n)
=e^{\tau}-\sum_{\rho}e^{\rho\tau}-\sum_{k\ge1}e^{-2k\tau}\qquad(\tau>0).
\]

\paragraph{A2 --- the reflection on the Hardy space.}
\(s\mapsto1-\bar s\) maps \(\operatorname{Re}s>\tfrac12\) outside itself.  It acts only on
boundary values, where it is complex conjugation; an \(H^2\) function with real boundary
values vanishes.  In the only well-defined reading the source's \(\Hplusplus\) is \(\{0\}\).
The antilinearity of \(J\) also makes \(\tfrac12(1\pm J)\) only real-linear.

\paragraph{A3 --- self-adjointness.}
A conjugation commuting with a symmetric operator gives equal deficiency indices, not
essential self-adjointness.  The operator \(-i(x\,d/dx+\tfrac12)\) is already self-adjoint on
\(L^2(\mathbb R^+,dx)\), with continuous spectrum \(\mathbb R\); self-adjointness is not the
bottleneck.

\paragraph{A4 --- discreteness and identification.}
The one-dimensional relation \([Q,P]=i\) does not produce a representation of an adelic Weyl
algebra; Stone--von Neumann for a locally compact abelian group yields \(L^2\) of the group,
not of \(\mathbb Q^+\backslash\mathbb A^1\).  The identification step uses Connes' trace formula
on the adele class space, whose validity is equivalent to RH; moreover the source's own formula
gives \(\mathcal L[\Sigma](s)=1/s-\log\zeta(s)+\mathcal L[R](s)\), not \(\log\zeta(s)\), and
\(\mathcal L[R]\) diverges.

\paragraph{A5 --- the set argument.}
The source assigns the same kernel \((\mathbb P,\mathrm{id},\log p)\) to \(\zeta\) and to every
self-dual automorphic \(L\)-function, and its operator stream is independent of the
\(L\)-function.  It would therefore assign one spectrum to all of them, while their zeros differ:
the first zero of \(\zeta\) is \(\tfrac12+14.1347251417\,i\), that of \(L(s,\chi_{-4})\) is
\(\tfrac12+6.0209489047\,i\).  The Davenport--Heilbronn function satisfies a Riemann-type
functional equation and has the zero \(0.8085171825+85.69934849\,i\) off the critical line.
Where the zeros lie is therefore not determined by the set \((\mathbb P,\log p)\) together with
the functional equation; it requires the multiplicative structure of the Euler product.

\paragraph{Theoremlet v4.8.0-r04.1.}
The Riemann source is not a proof source for RH.  None of the five guards of the release scope
depends on it, and the guard \(W_\gamma^{\rm glob}\), which forbids a global photon-clock/RH
theorem, is confirmed rather than affected.

\paragraph{Proof sketch.}
A1--A4 each break a link of the chain K1/L11 \(\to\) T1+T3 \(\to\) T4+T5 \(\to\) T6, and A5 shows
that the chain could not distinguish \(\zeta\) from other \(L\)-functions.  The guards concern
\(C5\), all-readout terminality, \(C6/T2\), \(W_\gamma^{\rm glob}\) and \(W_{\rm all}\); the
Companion already treated the RH channel as a local photon-clock probe, so no guarded statement
used RH as an input.  \(\square\)

\subsection*{v4.8.0-r04 standalone Heisenberg compensator}
\label{subsec:v480-r04-standalone-compensator}

\paragraph{Definition v4.8.0-r04.1.}
Let \(\mathcal H\) be a complex Hilbert space and \(J:\mathcal H\to\mathcal H\) a
\emph{linear} unitary involution, \(J^2=1\), \(J^*=J\).  The Heisenberg compensator and its
antisymmetric partner are
\[
C:=\tfrac12(1+J),\qquad A:=\tfrac12(1-J).
\]

\paragraph{Theoremlet v4.8.0-r04.2.}
\(C\) and \(A\) are complementary orthogonal projections, \(C+A=1\), \(CA=0\);
\(\operatorname{ran}C=\ker(J-1)\), \(\operatorname{ran}A=\ker(J+1)\); \(C(Jf)=C(f)\) and
\(A(Jf)=-A(f)\).  The terms \(\mathcal L_{\rm HC}\) and \(\mathcal L_{\rm frame}\) of
Theorem~\ref{thm:gauge-invariance} are \(J\)-invariant.  \(\mathcal L_{\rm kin}\) is
\(J\)-invariant if and only if \(JL_0J=L_0\).

\paragraph{Proof sketch.}
\(C^2=\tfrac14(1+2J+J^2)=C\) and \(C^*=C\) since \(J^*=J\); likewise for \(A\), and
\(CA=\tfrac14(1-J^2)=0\).  \(CJ=\tfrac12(J+1)=C\) and \(AJ=-A\).  \(\langle Jf,Jf\rangle=\langle
f,f\rangle\) by unitarity; \(\langle Jf,L_0Jf\rangle=\langle f,JL_0Jf\rangle\).  \(\square\)

\paragraph{Instances.}
(a) The \(e_1\)-reflection on \(\mathbb C\otimes\HH\) with the Euclidean inner product of the
coefficients is a linear unitary involution; on \(\HH\) it is an anti-automorphism.
(b) The linear reflection \(s\mapsto1-s\) restricted to the critical line acts on the ordinate as
the parity \(t\mapsto-t\), \((Jf)(t)=f(-t)\) on \(L^2(\mathbb R)\).  For (b), in logarithmic
coordinates \(L_0=-i(\partial_y+\tfrac12)\) and parity sends \(\partial_y\mapsto-\partial_y\), so
\(JL_0J\neq L_0\): the kinetic invariance is open for this instance.

The compensator therefore requires no RH statement.  What it loses relative to the earlier
text is the claim that its image is a space whose \(L_0\)-spectrum consists of zeta zeros.

\subsection*{v4.8.0-r04 standalone Sphaera cascade}
\label{subsec:v480-r04-standalone-cascade}

\paragraph{Definition v4.8.0-r04.2.}
Let \(\mathbb S\) be a fixed Hilbert space (a set, not a universal collection: ZF admits no
set of all sets), and let \(\mathbb S=S_N\supseteq S_{N-1}\supseteq\cdots\supseteq S_0\) be
closed subspaces with orthogonal projections \(\Pi_n:S_n\to S_{n-1}\).  The cascade is
\(\KC:=\Pi_1\circ\cdots\circ\Pi_N\).

\paragraph{Theoremlet v4.8.0-r04.3.}
\(\KC\) is the orthogonal projection of \(\mathbb S\) onto \(S_0\); in particular
\(\KC^2=\KC\), and \(\KC f=f\) if and only if \(f\in S_0\).

\paragraph{Proof sketch.}
For nested closed subspaces \(V\subseteq W\) the projections satisfy \(P_VP_W=P_V\); induction
over the chain gives \(\KC=P_{S_0}\).  \(\square\)

The Maximal Totality principle is read inside this definition: \(\KC\) is defined on
\(\mathbb S\) and maps every element of \(\mathbb S\) into \(S_0\).  It is not applied to
elements outside its domain.

\subsection*{v4.8.0-r04 repaired triole hierarchy}
\label{subsec:v480-r04-triole-hierarchy}

Use the Cayley--Dickson convention of the Riemann source,
\((p+q\ell)(r+s\ell)=(pr-\bar sq)+(sp+q\bar r)\ell\), with \(e_4=\ell\), \(e_5=e_1\ell\),
\(e_6=e_2\ell\), \(e_7=e_3\ell\).  Explicit multiplication gives the seven quaternionic triples
(lines of the Fano plane)
\[
(1,2,3),\ (1,4,5),\ (1,6,7),\ (2,4,6),\ (2,5,7),\ (3,4,7),\ (3,5,6).
\]

\paragraph{Theoremlet v4.8.0-r04.4.}
Three quaternionic triples cover all seven imaginary units of \(\OO\) if and only if they pass
through a common unit.  There are exactly seven such configurations, one per unit.  The source's
second triole \((e_4,e_5,e_6)\) is not quaternionic, since \(e_4e_5=e_1\).

\paragraph{Proof sketch.}
Two distinct lines of the Fano plane meet in exactly one point.  For three distinct lines,
\(|L_1\cup L_2\cup L_3|=9-|L_1\cap L_2|-|L_1\cap L_3|-|L_2\cap L_3|+|L_1\cap L_2\cap L_3|\),
which is \(7\) if the three lines are concurrent and \(6\) otherwise.  Each point lies on three
lines, so each point gives exactly one concurrent triple.  \(e_4e_5=\ell(e_1\ell)=e_1\) is the
source's own computation.  \(\square\)

\paragraph{Replacement.}
The triole configuration through \(e_4=\ell\) is
\[
\mathcal T_1'=(e_1,e_4,e_5),\qquad \mathcal T_2'=(e_2,e_4,e_6),\qquad
\mathcal T_3'=(e_3,e_4,e_7),
\]
which keeps the source's third triole and shares one common axis.  Equally admissible is the
configuration through \(e_1\), which keeps the inner triole \((e_1,e_2,e_3)\).

\paragraph{Reflection as algebra symmetry.}
The reflection negating \(e_1\) alone is not an automorphism of \(\OO\).  Of the
\(2^7\) sign changes exactly eight are automorphisms; those negating \(e_1\) negate four units
and fix one quaternionic line pointwise, for example \(\{e_1,e_2,e_4,e_7\}\) with fixed line
\((3,5,6)\).  A reflection meant as a symmetry of the algebra must have this form.

\paragraph{Above the octonions.}
In the sedenions the source's outer octo-triole \((f_8,\dots,f_{14})\) is not closed
(\(f_8f_9=f_1\)), and zero divisors exist, e.g. \((e_3+f_{10})(e_6-f_{15})=0\).  The hierarchy
continues algebraically, not geometrically; the Hopf structure ends at \(\OO\), as the source
itself states.  Theorem~\ref{thm:three-generations} uses \(e_4,e_5,e_6\) as an \(SU(3)\)
triplet, not as a quaternionic triole, and is not changed by this layer.

\subsection*{v4.8.0-r04 dependency map}
\label{subsec:v480-r04-dependency-map}

\begin{center}
\small
\begin{tabular}{@{}>{\raggedright\arraybackslash}p{0.31\textwidth}>{\raggedright\arraybackslash}p{0.30\textwidth}>{\raggedright\arraybackslash}p{0.33\textwidth}@{}}
\toprule
site & used from the Riemann source & status after r04\\
\midrule
Thm.~\ref{thm:unity}, equivalence clause & Stone--von Neumann (Lem.~13.2, Step~3) & open\\
Thm.~\ref{thm:unitary-equiv} & Stone--von Neumann, T4+T5 & open\\
Thms.~\ref{thm:bell-consistency}, \ref{thm:string-lqg-duality} & Stone--von Neumann & equivalence open\\
Thms.~\ref{thm:einstein}, \ref{thm:gue} & T6, \(\Sigma=1-\Theta+R\) & open; formula to be replaced (A1)\\
Thm.~\ref{thm:toe}, row ``Riemann Hypothesis'' & RH as a result & withdrawn: violates \(W_\gamma^{\rm glob}\)\\
Thm.~\ref{thm:gauge-invariance} & \(\mathrm{Adj}_J(L_0)=0\) & \(\mathcal L_{\rm HC}\), \(\mathcal L_{\rm frame}\) standalone; \(\mathcal L_{\rm kin}\) open\\
Thm.~\ref{thm:eom} & field equation & projector algebra kept; zero reading withdrawn\\
Thms.~\ref{thm:op2-existence}, \ref{thm:op2-selfadj} & T1+T3, T4+T5 & open (OP2 stays open)\\
Thm.~\ref{thm:partial-cancel}, Rem.~\ref{rem:explicit-formula-remainder} & explicit formula (Thm.~3.3) & to be replaced by A1\\
prime lattice \(\Gamma_P\) & classical arithmetic & unaffected\\
\S\ref{sec:rh} & RH as consequence & read as conjecture only\\
\bottomrule
\end{tabular}
\end{center}

\paragraph{Theoremlet v4.8.0-r04.5.}
After this layer the Heisenberg compensator, the Sphaera cascade and the triole hierarchy up to
\(\OO\) rest on Definitions v4.8.0-r04.1 and r04.2 and on Theoremlet r04.4, without any statement
of \cite{brendecke2026rh} beyond classical arithmetic.  The statements listed as open lose their
previous justification but are not refuted by this layer.

\paragraph{Proof sketch.}
Theoremlets r04.2--r04.4 use only Hilbert-space projections and explicit Cayley--Dickson
multiplication.  For the open items the audit A1--A5 removes the cited justification; whether
an independent justification exists is not decided here.  \(\square\)

\subsection*{v4.8.0-r04 ledger and guard update}
\label{subsec:v480-r04-ledger-guard}

\[
R_{\rm dec}^{v480\text{-}r04}=(0,0,0,1,1,1)
\]
for (standalone compensator; standalone cascade; repaired triole hierarchy; Stone--von Neumann
equivalence of the kernels open; trace-formula field equations open pending A1; kinetic
\(J\)-invariance open).

The layer spends no guard:
\[
R_{\rm glob}^{(v480\text{-}r04,+)}=(0,1,1,1,1),
\qquad
\Gamma_{\rm guard}^{v480\text{-}r04}=(1,1,1,1,1),
\qquad
\beta_4=1/256\ \text{unspent}.
\]
It does not imply exact finite \(C5\) zero closure, exact all-readout terminality,
\(C6/T2\), a global photon-clock/RH theorem, or global all-readout closure.

